\section{The local-information null space and the Type-II gate}
\label{sec:local-typeii-gate}

The preceding theorems preserve arbitrary labels that are already present
in a coefficient array.  They do not say whether a proposed label family
contains enough information to distinguish primes from balanced
semiprimes.  We make that question exact.

\subsection{An exact minimax formula for a local pushforward}

Let $I$ and $\Sigma$ be finite sets, let
$\sigma:I\to\Sigma$ be a signature map, and let $w=(w_i)_{i\in I}$ be
real weights.  The associated pushforward is
\begin{equation}
 (\cN_{\sigma,w}x)(s)
 =\sum_{\substack{i\in I\\\sigma(i)=s}}w_i x_i,
 \qquad s\in\Sigma.
 \label{eq:local-pushforward}
\end{equation}
For $v\in\R^I$, write $L_v(x)=\sum_i v_ix_i$.

\begin{theorem}[Exact local-data minimax radius]
\label{thm:local-minimax}
For $B>0$, even allowing an arbitrary nonlinear decoder
$\Phi:\ran\cN_{\sigma,w}\to\R$, one has
\begin{align}
 &\inf_{\Phi}
 \sup_{\|x\|_1\leq B}
 \left|L_v(x)-\Phi(\cN_{\sigma,w}x)\right|
 \notag\\
 &\qquad={}
 B\inf_{g:\Sigma\to\R}
 \max_{i\in I}|v_i-w_i g(\sigma(i))|.
 \label{eq:local-minimax}
\end{align}
Consequently exact recovery on the $\ell^1$ ball is possible if and only if
\begin{equation}
 v_i=w_i g(\sigma(i))
 \quad(i\in I)
 \label{eq:local-exact-recovery}
\end{equation}
for some $g$.  In the unweighted case $w_i=1$,
\begin{equation}
 e_B(v;\sigma)
 =\frac B2\max_{s\in\Sigma}
 \left(
 \max_{\sigma(i)=s}v_i-\min_{\sigma(i)=s}v_i
 \right).
 \label{eq:unweighted-fiber-oscillation}
\end{equation}
Empty fibers are omitted from the maximum.
\end{theorem}

\begin{proof}
Let $N=\cN_{\sigma,w}$.  Two points of the $\ell^1$ ball with the same
observation differ by an element of $\ker N$ of norm at most $2B$.
Conversely, if $z\in\ker N$ and $\|z\|_1\leq2B$, then $z/2$ and $-z/2$
are such a pair.  The smallest worst-case radius of all observation fibers
is therefore
\[
 B\sup_{\substack{z\in\ker N\\\|z\|_1\leq1}}|L_v(z)|.
\]
Finite-dimensional quotient duality identifies this norm with
\[
 B\,\dist_{\ell^\infty}
 \bigl(v,\ran N^*\bigr).
\]
Since $(N^*g)_i=w_i g(\sigma(i))$, this is
\eqref{eq:local-minimax}.  The zero-distance criterion gives
\eqref{eq:local-exact-recovery}.  When $w_i=1$, the best constant on each
fiber is the midpoint of the extreme $v_i$ values, proving
\eqref{eq:unweighted-fiber-oscillation}.
\end{proof}

Thus adding nonlinear post-processing does not repair a missing row.  The
obstruction is an exact null-space statement, not a limitation of a
particular learning rule or inversion algorithm.

\subsection{The square-root divisor-information dichotomy}

For $t\in I_X=(X,2X]\cap\N$ and $R\geq2$, define the complete local divisor
signature
\begin{equation}
 \sigma_R(t)=
 \bigl(\one_{d\mid t,\ d<t}\bigr)_{2\leq d\leq R}.
 \label{eq:divisor-signature}
\end{equation}
Thus only proper divisors are recorded, even if $R\geq t$.

\begin{proposition}[Square-root information gate]
\label{prop:square-root-information-gate}
The following statements hold.
\begin{enumerate}[label=\textup{(\roman*)}]
 \item If $R\geq\sqrt{2X}$, then $\sigma_R(t)=0$ for
 $t\in I_X$ if and only if $t$ is prime.
 \item Fix $0<\delta<1$.  For all sufficiently large $X$, uniformly for
 $R\leq(1-\delta)\sqrt{2X}$, there are a prime $p\in I_X$ and distinct
 primes $q_1,q_2>R$ such that
 \begin{equation}
  s=q_1q_2\in I_X,
  \qquad
  \sigma_R(p)=\sigma_R(s)=0.
  \label{eq:prime-semiprime-null-pair}
 \end{equation}
\end{enumerate}
\end{proposition}

\begin{proof}
If $t\leq2X$ is composite, it has a divisor $d$ with
$2\leq d\leq\sqrt t\leq\sqrt{2X}$, proving (i).

For (ii), put
\[
 c_0=\max\{1,(1-\delta)\sqrt2\}<\sqrt2
\]
and choose a closed interval $[c_1,c_2]$ strictly inside
$(c_0,\sqrt2)$.  The prime number theorem supplies two distinct primes in
$[c_1\sqrt X,c_2\sqrt X]$ for all sufficiently large $X$; after shrinking
the interval if necessary, every product of two points in it lies in
$(X,2X]$.  Call the primes $q_1,q_2$.  They exceed $R$, so their product
has zero signature.  A prime $p\in(X,2X]$, again supplied by the prime
number theorem, has the same signature.
\end{proof}

This is an information-resolution threshold.  It is not an analytic level
of distribution and does not assert that sieving to $\sqrt{2X}$ can be
carried out with a sufficiently accurate remainder.  It is a theorem about
the specified data interface $\sigma_R$: an algorithm supplied with the
integer $t$ itself or with additional arithmetic rows has more information.
Its relation to the classical parity limitation is therefore diagnostic,
not an alternative proof of that limitation \citep{Selberg1952Parity}.

Take the unweighted pushforward by $\sigma_R$ and view an unknown
$x\in\ell^1(I_X)$ as a signed mass distribution.  On the null pair in
\eqref{eq:prime-semiprime-null-pair}, the prime indicator takes values
$1$ and $0$, the Liouville function takes values $-1$ and $1$, and
$\Lambda$ takes values $\log p$ and $0$.  Hence
\cref{thm:local-minimax} gives, on the unit ball,
\begin{align}
 e_1(\one_{\mathbb P};\sigma_R)&=\frac12,
 \label{eq:prime-indicator-minimax}\\
 e_1(\lambda;\sigma_R)&=1,
 \label{eq:liouville-minimax}\\
 e_1(\Lambda;\sigma_R)&\geq\frac12\log p.
 \label{eq:lambda-minimax}
\end{align}
The first two equalities also use the global ranges
$\one_{\mathbb P}\in\{0,1\}$ and $\lambda\in\{-1,1\}$.

\subsection{A balanced factor row sees the null direction}

For coefficient sequences $\alpha,\beta$, define the bilinear observation
on a response vector $x$ by
\begin{equation}
 \cB_{\alpha,\beta}(x)
 =\sum_{m,n}\alpha_m\beta_n x(mn).
 \label{eq:typeii-row}
\end{equation}
For the pair in \eqref{eq:prime-semiprime-null-pair},
\begin{equation}
 \cB_{\delta_{q_1},\delta_{q_2}}(\delta_p)=0,
 \qquad
 \cB_{\delta_{q_1},\delta_{q_2}}(\delta_s)=1.
 \label{eq:typeii-separates-null-pair}
\end{equation}
Thus a single balanced factor row strictly enlarges this specified
divisor-signature data model: it separates two points that every test
$d\leq R$ identifies.

The same distinction is visible in the hard coefficient from
\eqref{eq:vaughan-identity}.  Remove the smooth weight and put
\begin{equation}
 H_{U,V}(N)=
 -\sum_{\substack{\ell k=N\\\ell>U,\ k>V}}
 \Lambda(\ell)\beta_V(k).
 \label{eq:finite-vaughan-hard-row}
\end{equation}
If $p,q_1,q_2>\max\{U,V\}$ and $s=q_1q_2$, then
\begin{equation}
 H_{U,V}(p)=0,
 \qquad
 H_{U,V}(s)=-\log s.
 \label{eq:vaughan-hard-null-witness}
\end{equation}
Indeed $\beta_V(q_i)=1$.  In Vaughan's identity the $d=1$ logarithmic
packet contributes $+\log N$ to both inputs; the hard packet cancels it
exactly for the semiprime and leaves it untouched for the prime.

Equations \eqref{eq:typeii-separates-null-pair} and
\eqref{eq:vaughan-hard-null-witness} certify that this balanced bilinear
row contains information absent from the specified local signature.  They
provide no cancellation estimate for a long Type-II sum.  That analytic
estimate is the gate that remains.
